\begin{resumo}
O presente documento consolida o relatório formal de Busca de Anterioridade e Prospecção Tecnológica desenvolvido para a concepção, proteção intelectual e validação do sistema SIGEA-GTT (Sistema Inteligente de Gestão de Eventos Adversos -- Global Trigger Tool). A plataforma foi projetada como uma solução computacional inovadora para identificação ativa, rastreamento sistemático e auditoria de riscos clínicos baseada na metodologia \textit{Global Trigger Tool} (GTT) do \textit{Institute for Healthcare Improvement} (IHI), incorporando como diferencial inventivo a sua aplicação pedagógica síncrona e não-punitiva na formação de graduandos em Enfermagem e Medicina. A prospecção do estado da técnica foi realizada em 02 de setembro de 2026 em cinco bases referenciais globais e nacionais de patentes e artigos científicos revisados por pares: WIPO Patentscope, Scopus, Web of Science, SciELO e Portal de Periódicos CAPES, operando com formulação booleana tripartite (Segurança do Paciente, Engenharia de Software e Educação/Treinamento). Na base WIPO Patentscope, o refinamento na classe IPC G16H revelou 86 documentos correlatos, dentre os quais quatro patentes destacam-se como estado da técnica mais próximo (CN118888101, CN111627533, CN113379384 e US20210074398), todas voltadas estritamente a fluxos administrativos corporativos hospitalares de faturamento ou sequenciamento curricular estático. Na literatura científica, três artigos delimitam a fronteira do conhecimento em validação interdisciplinar, usabilidade psicométrica e mitigação do fenômeno da ``segunda vítima''. A análise técnica comprova que o SIGEA-GTT atende integralmente aos requisitos de Novidade, Atividade Inventiva e Aplicação Industrial preconizados pela Lei de Propriedade Industrial (Lei n.º 9.279/1996), fundamentando o avanço para o depósito de patente de invenção e registro de programa de computador junto ao Instituto Nacional da Propriedade Industrial (INPI).

\textbf{Palavras-chave}: Busca de Anterioridade. Prospecção Tecnológica. Patentes. Global Trigger Tool. Eventos Adversos. Segurança do Paciente. Engenharia de Software. SIGEA-GTT.
\end{resumo}
